\section{Penalaran Umum}
\subsection{Logika Proposional (Hubungan Sebab-Akibat)}
Di dalam logika proposional biner, kalimat bahasa indonesia direduksi menjadi variabel proposisi $p, q, r \dots \in \{B, S\}$.

\subsubsection{Impilkasi Kondisional}
Bentuk ini menyatakan hubungan syarat (anteseden) dan konsekuensei (konsekuen).
\begin{itemize}
\item \textbf{Ekstraksi Kalimat: } ``Jika $p$ maka $q$'', ``$q$ saat $p$'', ``$q$ apabila $p$'', ``Setiap kali $p$ terjadi, maka $q$ terjadi''.
\item \textbf{Notasi Formal: } \[p \implies q \]
\item \textbf{Ekuivalensi Tautologis (Hukum Kontraposisi \& Disjungsi: )} \[p \implies q \equiv \neg q \implies \neg p \equiv \neg p \lor q \]
\end{itemize}

\subsubsection{Konjungsi}
Bentuk ini menyatakan dua peristiwa yang wajib terjadi bersamaan agar nilai kebenaran bernilai benar (B).
\begin{itemize}
\item \textbf{Ekstraksi Kalimat: } ``$p$ dan $q$'', ``$p$ tetapi $q$'', ``$p$ walaupun $q$''.
\item \textbf{Notasi Formal: } \[p \land q\]
  \item \textbf{Hukum De Morgan (Negasi Konjungsi: )} \[ \neg (p \land q) \equiv \neg p \lor \neg q \]
\end{itemize}

\subsubsection{Disjungsi Inklusif}
Bentuk ini menyatakan setidaknya salah satu dari kedua peristiwa terjadi.
\begin{itemize}
\item \textbf{Ekstraksi Kalimat: } ``$p$ atau $q$''.
\item \textbf{Notasi Formal: } \[p \lor q\]
\item \textbf{Hukum De Morgan (Negasi Disjungsi:) } \[ \neg (p \lor q) \equiv \neg p \land \neg q \]
\end{itemize}

\subsection{Logika Predikat dan Kuantor (Semua vs Sebagian)}
Untuk tipe soal yang melibatkan batasan jumlah objek, gunakan pendekatan kuantor formal yang berkorelasi langsung dengan Aljabar Himpunan (Diagram Venn).

\subsubsection{Kuantor Universal (Universal Quantifier)}
Menyatakan seluruh semesta pembicaraan terikat pada sifat tertentu. Secara geometris merupakan hubungan \textbf{Himpunan Bagian (Subset)}.
\begin{itemize}
\item \textbf{Ekstraksi Kalimat: } ``Semua $A$ adalah $B$'', ``Setiap $A$ memiliki sifat $B$''.
\item \textbf{Notasi Logika Predikat: } \[ \forall x \, (A(x) \implies B(x))\]
\item \textbf{Representasi Aljabar Himpunan: } \[ A \subseteq B \quad \iff \quad A \cap B = A\]
\item \textbf{Aturan Negasi: }
  \[ \neg \big[ \forall x \, (A(x) \implies B(x)) \big] \equiv \exists x \, (A(x) \land \neg B(x)) \]
  \textit{Dibaca: ``Ada/Sebagian A yang Bukan B''}
\end{itemize}

\subsubsection{Kuantor Eksistensial (Existensial Quantifier)}
Menyatakan setidaknya ada minimal satu objek dalam semesta pembicaraan yang memenuhi syarat. Secara geometris merupakan hubungan \textbf{Irisan (Intersection)}.
\begin{itemize}
\item \textbf{Ekstraksi Kalimat: } ``Sebagian $A$ adalah $B$'', ``Beberapa $A$ memiliki sifat $B$'', ``Ada $A$ yang termasuk $B$'', ``Sementara $A$ adalah $B$''.
\item \textbf{Notasi Logika Predikat: } \[ \exists x \, (A(x) \land B(x)) \]
\item \textbf{Representasi Aljabar Himpunan: } \[ A \cap B \neq \emptyset \]
\item \textbf{Aturan Negasi: }\[ \neg \big[ \exists x \, (A(x) \land B(x))\big] \equiv \forall x \, (A(x) \implies \neg B(x)) \]
  \textit{(Dibaca: ``Semua A bukan B'')}
\end{itemize}

\subsection{Aturan Inferensi (Mesin Penarikan Kesimpulan)}
Gunakan skema deduktif formal di bawah ini. Jika premis soal dapat disusun menjadi bentuk di atas garis, make kesimpulan di bawah tanda $\therefore$ adalah jawaban mutlak yang valid.

\subsubsection{Modus Ponens}
\[ \frac{p \implies q, \quad p}{\therefore q}\]

\subsubsection{Modus Tollens}
\[ \frac{p \implies q, \quad q \implies r}{\therefore \neg p}\]

\subsubsection{Silogisme Hipotetis (Rantai Transitif Proposisi)}
\[ \frac{p \implies q, \quad q \implies r}{\therefore p \implies r \equiv \neg r \implies \neg p}\]

\subsubsection{Silogisme Disjungtif}
\[ \frac{p \lor q, \quad \neg p}{\therefore q} \qquad \text{dan} \qquad \frac{p \lor q, \quad \neg q}{\therefore p} \]

\subsubsection{Pola Divergen (Satu Anteseden, Dua Konsekuen)}
\[ \frac{p \implies q, \quad p \implies r}{\therefore p \implies (q land r)} \]

\subsubsection{Silogisme Kuantor Transitif (Semua + Semua)}
\[ \frac{\forall x \, (A(x) \implies B(x)), \quad \forall x \, (B(x) \implies C(x))}{\therefore \forall x \, (A(x) \implies C(x)) \quad \iff \quad A \subseteq C} \]
\textit{(Dibaca: ``Semua A adalah C'')}

\subsubsection{Silogisme Kuantor Partikular (Semua + Sebagian)}
\[ \frac{\forall x \, (A(x) \implies B(x)), \quad \exists x \, (C(x) \land A(x))}{\therefore \exists x \, (C(x) \land B(x)) \quad \iff \quad C \cap B \neq \emptyset} \]
\textit{(Dibaca: ``Sebagian C Adalah B'')}

\subsection{Tabel Fallacy}
Opsi pengecoh di UTBK selalu dibuat menggunakan pola inferensi yang cacat secara formal (\textit{invalid argument}). Jangan pernah memilih opsi yang memiliki pola berikut:

\begin{table}[h]
  \centering
  \begin{tabular}{lcc}
    \toprule
    \textbf{Nama Fallacy} & \textbf{Susunan Premis} & \textbf{Kesimpulan Salah pada Opsi} \\
    \midrule
    Affirming the Consequent & $p \implies q, \quad q$ & $\therefore p$ \\
    \addlinespace
    Denying the Antecedent & $p \implies q, \quad \neg p$ & $\therefore \neg q$ \\
    \addlinespace
    Kuantor Eksistensial Ganda & $A \cap B \neq \emptyset, \quad B \cap C \neq \emptyset$ & $\therefore A \cap C \neq \emptyset$ \\
    \addlinespace
    Konversi Kuantor Universal & $A \subseteq B$ & $\therefore B \subseteq A$ \\
    \bottomrule
  \end{tabular}
\end{table}
